\documentclass{zhvt-classic}

% 使用本地源流明体（GenRyuMin2 TW）按文件路径加载，不依赖系统字体名
% 版框（内框／界栏／书口侧栏／鱼尾）改用与版心印章同一印泥棕色：
%   印章图片 yshs.png 实色 = RGB(135,68,52) = HTML 874434
\usepackage{xcolor}
\usepackage{tikz}
% sealbrown 双注册：xcolor 供 TikZ/pgf，L3 供 cls 版框
\definecolor{sealbrown}{HTML}{874434}
\ExplSyntaxOn
\color_set:nnn { sealbrown } { HTML } { 874434 }
\ExplSyntaxOff

\zhvtset{
  grid color       = sealbrown,
  main font        = GenRyuMin2TW-SB,
  main font path   = C:/Users/Administrator/texmf/fonts/opentype/public/GenRyuMin2TW/,
jiazhu font      = GenRyuMin2TW-M,
  jiazhu font path = C:/Users/Administrator/texmf/fonts/opentype/public/GenRyuMin2TW/,
  % 页码框容纳 4 字：正常页「九十七」+前置空 1 格=4 行恰好；
  % 100+ 页「一百零一」共 5 行超出，末字「一/二/三」被框裁掉（2026-10-04 实测），故放宽到 5
  shukou page number chars = 5,
  % 书口书名框高从默认 6 字减到 4 字，把省出的 2 字高拨给 chapter 框，
  % 使 6 字书名+卷名放进 9 字高的 chapter 框后不再压线（2026-10-05）
  shukou title chars = 4,
}

\title{史记}[穰侯列传]
\author{汉·司马迁}
\maker{闲叟漫记·重製}


% ===== 索隐述赞样式（2026-11-08 四次定稿 · 愚大人亲手改定参数）=====
% ⚠ 竖排前提：自定义字体家族必须带 RawFeature={vertical:+vert}；
%   simhei.ttf 加 +vhal 会整体加载失败静默回退，只写 +vert；simfang 同理。
\newCJKfontfamily{\shuzanfont}[Path=C:/Windows/Fonts/,RawFeature={vertical:+vert}]{simhei.ttf}
\newCJKfontfamily{\shuzanbody}[Path=C:/Windows/Fonts/,RawFeature={vertical:+vert}]{simfang.ttf}
% 题头：152x30pt 圆角 6pt 棕底白字；baseline (0,2) 使整框落入界行格子不压版框线；
% 白字节点 x 不补偿、y +8.2pt（像素级墨迹实测 dx/dy≈0，L/R 留白约 20pt）。
\newcommand{\shuzanhead}[1]{%
  \begin{tikzpicture}[baseline={(0,2)}]
    \fill[sealbrown, rounded corners=6pt]
      (-76pt,-15pt) rectangle (76pt,15pt);
    \node[text=white, inner sep=0pt, anchor=center] at (0pt,8.2pt)
      {\shuzanfont\fontsize{23}{23}\selectfont #1};
  \end{tikzpicture}%
}
% ★ 修正书口卷名位置：chapter 框起始位置改回 cls 默认值，避免全局影响。
%   只在每个卷首页通过 \juan 命令里的空命令微调该页位置（安全，不会跨页污染）。
\ExplSyntaxOn
\AtBeginDocument{
  \dim_set:Nn \l__zhvt_shukou_chapter_title_start_dim
    { \l__zhvt_shukou_title_height_dim + \l__zhvt_font_size_dim * 2 }
  \dim_set:Nn \l__zhvt_shukou_page_number_start_dim
    { \l__zhvt_shukou_chapter_title_start_dim + \l__zhvt_shukou_chapter_title_height_dim }
}
\ExplSyntaxOff

\begin{document}

% ★★ 版心（书口）排布覆写：把书名放到鱼尾之下（2026-10-04）
%
%   cls 默认版心自上而下：书名（title 静态框，在鱼尾之上）→ 鱼尾 → 篇名（chapter 动态框）→ 页码
%   本模板改为：          鱼尾（位置不动）→ 书名 → 篇名 → 页码分割线 → 页码
%
%   手法：把 cls 放在 title 静态框里的书名，改放进 chapter 动态框（鱼尾之下）；
%         title 静态框留空。书名取 \l__zhvt_title_str，即 \title{...} 的第一参数，
%         所以换书时只改 \title，不必在此重复写书名。
%
%   ⚠ 三条必须遵守的写法（2026-10-04 实测）：
%     ① 汉字必须用 \hbox{...} 包住。flowfram 会把框内容装进一个宽 14.4mm（= 一个 line_sep
%        = 书口栏宽）的纵向 minipage 并按该宽度折行，而整页靠 PDF /Rotate 90 显示、
%        框内文字本身不变向 —— 不包 \hbox 时「围炉夜话」4 字会被逐字折成 4 行，
%        在版面上就成了横排 4 字并横向溢出版框。
%     ② 竖排下**一个 \hbox 就是一整列**。书名与篇名要同处一列，必须放进同一个 \hbox；
%        分成两个 \hbox 时，后一个会被当成新的一列、甩到隔壁栏去
%        （实测「序」跑到书口栏外侧 14.4mm 的正文列上）。中间的 \quad 在框内是横向空白，
%        显示出来正好是书名与篇名之间的纵向间隔。
%     ③ ★ 内容盒要写成 **与框等宽的定宽盒** `\hbox to \l__zhvt_shukou_chapter_title_height_dim{…\hfil}`。
%        实测：chapter 框是把内容盒沿纵向**居中**摆放的，内容宽度一变，书名就跟着漂：
%        自然宽度的 \hbox（书名 38.4mm）落在 x=157.2mm，加篇名后（57.6mm）落在 x=152.4mm，
%        目录页（无篇名）与正文页的书名就差半格。做成等宽盒后，无论有没有篇名，
%        内容固定从框顶起排，实测两页均落在 x=142.8mm，不再漂。
\ExplSyntaxOn
% 把书口 chapter 框（书名+篇名）起始位置下移 1 字高，避免 4 字卷名顶到鱼尾
% ⚠ cls 的 \__zhvt_calculate_dimensions: 在 \AtBeginDocument 里执行，会重算该 dim；
%   必须把覆写挂到 begin document 之后，否则会被二次覆盖（2026-10-05 实测）。
% (旧方案已移到导言区，直接修改 chapter 框起始位置)
% 卷标题命令 #1=书口篇名(鱼尾) #2=正文卷名（含卷序，卷名内空格用 ~）
%
% ★★ 目录页两项修订（2026-10-05，愚大人指出）★★
%   ① 页码起点统一对齐：
%      原状 —— 目录项写成 \par\hspace{4em}卷名\hspace{4em}页码，页码起点 = 卷名实宽 + 4em。
%      卷名实宽随字数变：「卷一~梁惠王上」6 字 +~、「卷七~离娄上」5 字 +~，
%      差 1 个字位 → 卷七~卷十 4 条的页码整体比其余 10 条高 27.21pt（一个字），看着上下不齐。
%      改法 —— 卷名补全角空格（宽 = 1 字 = 27.215pt）补足到 **7 字符位**再写进目录，
%      14 条页码起点遂全部落在同一条竖线上（起点对齐；末字随页码字数自然参差）。
%      ⚠ 只对目录里的副本 \l_tmpc_tl 补，正文卷名仍用 \l_tmpa_tl 原样 —— 正文像素不受影响。
%   ② 卷名整体上移半字位：前导 \hspace 由 4em 改 3.5em。
%      实测原状目录块在版框内偏下（天头留白 139.6pt / 地脚留白 110.7pt），
%      上移 0.5em = 13.6pt 后上下基本均分，视觉居中。
\NewDocumentCommand { \juan } { mm } {
	\clearpage
	\tl_set:Nn \l__zhvt_chapter_title_tl { #1 }
	\tl_set:Nn \l_tmpa_tl { #2 }
	\tl_set_eq:NN \l_tmpc_tl \l_tmpa_tl
	\int_step_inline:nn { 7 - \str_count:n { #2 } }
	  { \tl_put_right:Nn \l_tmpc_tl { 　 } }
	\tl_set:Nx \l_tmpb_tl { \zhnum { page } }
	\tl_put_left:Nn \l_tmpb_tl { \hspace { 4em } }
	\tl_put_left:Nn \l_tmpb_tl { \l_tmpc_tl }
	\tl_put_left:Nn \l_tmpb_tl { \par \hspace { 3.5em } }
	\addtocontents { toc } { \l_tmpb_tl }
	\skip_horizontal:n { \l__zhvt_font_size_dim * 4 }
	\l_tmpa_tl
	\par
	\thispagestyle { empty }
}
\cs_set_protected:Npn \__zhvt_set_shukou_content: {
	% ★ 鱼尾上方留白（title 静态框）嵌入闲叟钤印：
	%   整页 /Rotate 90，故此处把图用 \rotatebox{90} 转正为竖章；
	%   尺寸卡在版心条幅内（书内约 7mm 宽 × 33.8mm 高，留余量不压鱼尾）。
	\setstaticcontents * { title } {
		\vbox_to_ht:nn { \l__zhvt_baseline_skip_dim } {
			\vskip 14.2pt
			\hbox_to_wd:nn { \l__zhvt_shukou_title_height_dim } {
				\hfil
				\rotatebox [ origin = c ] { 90 } { \includegraphics [ width = 7mm , height = 33.8mm ] { yshs.png } }
				\hfil
			}
		}
	}
% ★ 版心（书口）统一改法（2026-11-16 愚大人定）：18pt＋框内居中＋书口文字同版框色
%   ① 原 27.2pt 时 7 字长篇名（如「廉颇蔺相如列传」）末字溢出 chapter 框、压住页码；
%      改 \fontsize{18pt}{18pt}，前后 \hfil 令字串在框内居中（不随篇名字数漂移）。
%   ② \vbox{\vskip 8.2pt} 定量下移，抵消缩字号后栏宽方向 8.2pt 上偏（像素实测 y 中心 543.3 = 原版）。
%      ⚠ \raisebox / \vspace* 在 flowfram 框内无效；换字号须同步换补偿：16pt→10.1、19pt→7.4。
%   ③ chapter 与 page~number 均置 \color{sealbrown}（HTML 874434），与版框／界行同色；正文仍黑。
%   ⚠ 页码字位、钤印、鱼尾、界行、正文一律未动；五部合订本同此口径。
	\setdynamiccontents * { chapter }
	{ \color{sealbrown} \vbox { \vskip 8.2pt \hbox to \l__zhvt_shukou_chapter_title_height_dim
		{ \hfil \fontsize{18pt}{18pt}\selectfont \l__zhvt_title_str \quad \l__zhvt_chapter_title_tl \hfil } } }
	\setdynamiccontents * { page~number } { \color{sealbrown} \quad \zhnum { page } }
}
\__zhvt_set_shukou_content:
\ExplSyntaxOff

\maketitle

% 单篇，无多卷目录，不渲染目录页

\juan{穰侯列传}{卷七十二~穰侯列传第十二}
\hspace*{2em}穰侯魏冉者。秦昭王母宣太后弟也。其先楚人。姓羋氏。\\
\hspace*{2em}秦武王卒。无子。立其弟为昭王。昭王母故号为羋八子。及昭王即位。羋八子号为宣太后。宣太后非武王母。武王母号曰惠文后。先武王死。宣太后二弟。其异父长弟曰穰侯。姓魏氏。名冉。同父弟曰羋戎。为华阳君。而昭王同母弟曰高陵君。泾阳君。而魏冉最贤。自惠王。武王时任职用事。武王卒。诸弟争立。唯魏冉力为能立昭王。昭王即位。以冉为将军。卫咸阳。诛季君之乱。而逐武王后出之魏。昭王诸兄弟不善者皆灭之。威振秦国。昭王少。宣太后自治。任魏冉为政。\\
\hspace*{2em}昭王七年。樗里子死。而使泾阳君质於齐。赵人楼缓来相秦。赵不利。乃使仇液之秦。请以魏冉为秦相。仇液将行。其客宋公谓液曰。﹁秦不听公。楼缓必怨公。公不若谓楼缓曰﹃请为公毋急秦﹄。秦王见赵请相魏冉之不急。且不听公。公言而事不成。以德楼子。事成。魏冉故德公矣。﹂於是仇液从之。而秦果免楼缓而魏冉相秦。\\
\hspace*{2em}欲诛吕礼。礼出奔齐。昭王十四年。魏冉举白起。使代向寿将而攻韩。魏。败之伊阙。斩首二十四万。虏魏将公孙喜。明年。又取楚之宛。叶。魏冉谢病免相。以客卿寿烛为相。其明年。烛免。复相冉。乃封魏冉於穰。复益封陶。号曰穰侯。\\
\hspace*{2em}穰侯封四岁。为秦将攻魏。魏献河东方四百里。拔魏之河内。取城大小六十馀。昭王十九年。秦称西帝。齐称东帝。月馀。吕礼来。而齐。秦各复归帝为王。魏冉复相秦。六岁而免。免二岁。复相秦。四岁。而使白起拔楚之郢。秦置南郡。乃封白起为武安君。白起者。穰侯之所任举也。相善。於是穰侯之富。富於王室。\\
\hspace*{2em}昭王三十二年。穰侯为相国。将兵攻魏。走芒卯。入北宅。遂围大梁。梁大夫须贾说穰侯曰。﹁臣闻魏之长吏谓魏王曰。﹃昔梁惠王伐赵。战胜三梁。拔邯郸。赵氏不割。而邯郸复归。齐人攻卫。拔故国。杀子良。卫人不割。而故地复反。卫。赵之所以国全兵劲而地不并於诸侯者。以其能忍难而重出地也。宋。中山数伐割地。而国随以亡。臣以为卫。赵可法。而宋。中山可为戒也。秦。贪戾之国也。而毋亲。蚕食魏氏。又尽晋国。战胜暴子。割八县。地未毕入。兵复出矣。夫秦何厌之有哉。今又走芒卯。入北宅。此非敢攻梁也。且劫王以求多割地。王必勿听也。今王背楚。赵而讲秦。楚。赵怒而去王。与王争事秦。秦必受之。秦挟楚。赵之兵以复攻梁。则国求无亡不可得也。原王之必无讲也。王若欲讲。少割而有质。不然。必见欺。﹄此臣之所闻於魏也。原君之以是虑事也。周书曰﹃惟命不于常﹄。此言幸之不可数也。夫战胜暴子。割八县。此非兵力之精也。又非计之工也。天幸为多矣。今又走芒卯。入北宅。以攻大梁。是以天幸自为常也。智者不然。臣闻魏氏悉其百县胜甲以上戍大梁。臣以为不下三十万。以三十万之众守梁七仞之城。臣以为汤。武复生。不易攻也。夫轻背楚。赵之兵。陵七仞之城。战三十万之众。而志必举之。臣以为自天地始分以至于今。未尝有者也。攻而不拔。秦兵必罢。陶邑必亡。则前功必弃矣。今魏氏方疑。可以少割收也。原君逮楚。赵之兵未至於梁。亟以少割收魏。魏方疑而得以少割为利。必欲之。则君得所欲矣。楚。赵怒於魏之先己也。必争事秦。从以此散。而君後择焉。且君之得地岂必以兵哉。割晋国。秦兵不攻。而魏必效绛安邑。又为陶开两道。几尽故宋。卫必效单父。秦兵可全。而君制之。何索而不得。何为而不成。原君熟虑之而无行危。﹂穰侯曰。﹁善。﹂乃罢梁围。\\
\hspace*{2em}明年。魏背秦。与齐从亲。秦使穰侯伐魏。斩首四万。走魏将暴鸢。得魏三县。穰侯益封。\\
\hspace*{2em}明年。穰侯与白起客卿胡阳复攻赵。韩。魏。破芒卯於华阳下。斩首十万。取魏之卷。蔡阳。长社。赵氏观津。且与赵观津。益赵以兵。伐齐。齐襄王惧。使苏代为齐阴遗穰侯书曰。﹁臣闻往来者言曰﹃秦将益赵甲四万以伐齐﹄。臣窃必之敝邑之王曰﹃秦王明而熟於计。穰侯智而习於事。必不益赵甲四万以伐齐﹄。是何也。夫三晋之相与也。秦之深雠也。百相背也。百相欺也。不为不信。不为无行。今破齐以肥赵。赵。秦之深雠。不利於秦。此一也。秦之谋者。必曰﹃破齐。弊晋。楚。而後制晋。楚之胜﹄。夫齐。罢国也。以天下攻齐。如以千钧之弩决溃筴也。必死。安能弊晋。楚。此二也。秦少出兵。则晋。楚不信也。多出兵。则晋。楚为制於秦。齐恐。不走秦。必走晋。楚。此三也。秦割齐以啖晋。楚。晋。楚案之以兵。秦反受敌。此四也。是晋。楚以秦谋齐。以齐谋秦也。何晋。楚之智而秦。齐之愚。此五也。故得安邑以善事之。亦必无患矣。秦有安邑。韩氏必无上党矣。取天下之肠胃。与出兵而惧其不反也。孰利。臣故曰秦王明而熟於计。穰侯智而习於事。必不益赵甲四万以代齐矣。﹂於是穰侯不行。引兵而归。\\
\hspace*{2em}昭王三十六年。相国穰侯言客卿灶。欲伐齐取刚。寿。以广其陶邑。於是魏人范睢自谓张禄先生。讥穰侯之伐齐。乃越三晋以攻齐也。以此时奸说秦昭王。昭王於是用范睢。范睢言宣太后专制。穰侯擅权於诸侯。泾阳君。高陵君之属太侈。富於王室。於是秦昭王悟。乃免相国。令泾阳之属皆出关。就封邑。穰侯出关。辎车千乘有馀。\\
\hspace*{2em}穰侯卒於陶。而因葬焉。秦复收陶为郡。\\
\hspace*{2em}太史公曰。穰侯。昭王亲舅也。而秦所以东益地。弱诸侯。尝称帝於天下。天下皆西乡稽首者。穰侯之功也。及其贵极富溢。一夫开说。身折势夺而以忧死。况於羁旅之臣乎。\\

% ===== 索隐述赞（司马贞《史记索隐》）· 统一样式 =====
% 题头「索隐述赞曰」五字并入反白框；正文每四句一列，定宽盒 \hfil 居列中；
% \fontsize{24}{41} 列距对齐全书界行（121+41k）；\clearpage 使述赞独占一页。
\clearpage
\shuzanhead{索隐述赞曰}\linebreak
{\shuzanbody\fontsize{24}{41}\selectfont
\hbox to \hsize{\hfil 穰侯智识。应变无方。内倚太后。外辅昭王。\hfil}\linebreak
\hbox to \hsize{\hfil 四登相位。再列封疆。摧齐挠楚。破魏围梁。\hfil}\linebreak
\hbox to \hsize{\hfil 一夫开说。忧愤而亡。\hfil}%
\par}
\end{document}
